% =============================================================
% บทที่ 5 บทสรุปและข้อเสนอแนะ
% =============================================================

% -----------------------------------------------------------
\section{สรุปผลการดำเนินโครงงาน}
% -----------------------------------------------------------

โครงงานได้พัฒนาระบบรวบรวมและสังเคราะห์ข่าวภาษาไทยตั้งแต่การดึงข่าวจากสำนักข่าว 14 แห่ง การสกัดและจำแนกข่าว 9 หมวดหมู่ ไปจนถึงการสรุป วิเคราะห์ และนำเสนอข่าวพร้อมจัดกลุ่มประเด็นและจัดอันดับกระแสข่าว ผลประเมินด้านการจำแนกข่าวได้ Macro F1 เท่ากับ 0.8749 และความถูกต้องร้อยละ 87.78 ส่วนการจัดกลุ่มประเด็นข่าวได้ B-cubed F1 เท่ากับ 0.8026 การทดสอบการยอมรับของผู้ใช้ผ่านครบ 15 กรณี และชุดทดสอบอัตโนมัติผ่าน 423 กรณี นอกจากนี้ ผู้ใช้ 9 คนประเมินความพึงพอใจโดยรวมในระดับดีมาก โดยมีค่าเฉลี่ยตามผลสรุปในตารางเท่ากับ 4.86 จาก 5 และคะแนนแนวโน้มการแนะนำบอกต่อเฉลี่ย 6.33 จาก 7

\section{ข้อจำกัดของระบบ}
% -----------------------------------------------------------

จากการพัฒนาและทดสอบระบบจริง คณะผู้จัดทำได้สรุปข้อจำกัดของระบบออกเป็น 2 ด้านหลัก ดังนี้

\subsection{ข้อจำกัดด้านการทำงาน}

\begin{enumerate}
    \item \textbf{การเปลี่ยนแปลงโครงสร้างหน้าเว็บต้นทาง (DOM Mutation)}\\*
    เว็บไซต์ของสำนักข่าวออนไลน์อาจมีการปรับปรุงรูปแบบและโครงสร้างเอกสาร HTML หรือ CSS Selector ซึ่งอาจส่งผลให้ตัวแยกวิเคราะห์ (Parser) สกัดเนื้อหาบทความได้ไม่ครบถ้วน จนกว่าจะได้รับการปรับปรุงการตั้งค่าตัวสกัดให้ตรงกับโครงสร้างใหม่
    \item \textbf{การพึ่งพาบริการโมเดลภาษาขนาดใหญ่ภายนอก (External LLM Dependency)}\\*
    การสังเคราะห์เนื้อหาข่าวและการวิเคราะห์อารมณ์พึ่งพาการเชื่อมต่อ KKU Gen.AI API ซึ่งในกรณีที่เครือข่ายภายนอกขัดข้องหรือโควตาการเรียกใช้งานต่อนาทีเต็ม ระบบจะต้องสลับไปแสดงผลข้อความสรุปสำรองเชิงฮิวริสติก (Heuristic Fallback) โดยอัตโนมัติ
    \item \textbf{ขอบเขตของโมเดลการรู้จำชื่อเฉพาะ (NER Scope)}\\*
    โมเดลสกัดเอนทิตีภาษาไทย \texttt{thainer-v2} ถูกฝึกสอนจากคลังข้อมูลในอดีต ทำให้กรณีที่มีชื่อบุคคล สถานที่ หรือศัพท์สแลงเฉพาะกิจเกิดใหม่ในกระแสข่าวปัจจุบัน โมเดลอาจยังไม่สามารถรู้จำชื่อเฉพาะเหล่านั้นได้อย่างครบถ้วน
    \item \textbf{การประมวลผลเฉพาะข่าวภาษาไทย}\\*
    ระบบได้รับการออกแบบและฝึกสอนโมเดลเพื่อรองรับการประมวลผลภาษาไทยเป็นหลัก จึงยังไม่สามารถรวมกลุ่มบทความข่าวภาษาต่างประเทศเข้ากับประเด็นข่าวภาษาไทยในเหตุการณ์เดียวกันได้
\end{enumerate}

\subsection{ข้อจำกัดด้านฮาร์ดแวร์และสภาพแวดล้อม}

\begin{enumerate}
    \item \textbf{การใช้ทรัพยากรของ Playwright Headless Browser}\\*
    การเปิดใช้อินสแตนซ์ของเบราว์เซอร์ Chromium เพื่อสกัดหน้าเว็บที่เรนเดอร์ด้วย JavaScript และผ่านด่านตรวจบอต จำเป็นต้องใช้หน่วยความจำหลัก RAM และรอบเวลาประมวลผลของ CPU สูงกว่าคำขอ HTTP ปกติ ทำให้ระบบต้องจำกัดจำนวนการเปิดหน้าต่างเบราว์เซอร์พร้อมกันเพื่อรักษาเสถียรภาพของเครื่องแม่ข่าย
    \item \textbf{การประมวลผลโมเดล Transformer บนหน่วยประมวลผลกลาง (CPU-only)}\\*
    บนเครื่องแม่ข่ายคลาวด์ที่ไม่มีการติดตั้งหน่วยประมวลผลกราฟิก (GPU) การสร้างเวกเตอร์ความหมาย 768 มิติด้วยโมเดล WangchanBERTa จะใช้เวลาประมวลผลเฉลี่ยประมาณ 48 มิลลิวินาทีต่อข้อความ ซึ่งอาจส่งผลให้รอบการจัดกลุ่มข่าวใช้เวลาเพิ่มขึ้นเมื่อมีข่าวเข้าสู่ระบบพร้อมกันจำนวนมาก
    \item \textbf{การแข่งขันในการเขียนข้อมูลของ SQLite}\\*
    แม้ว่าฐานข้อมูล SQLite จะเปิดใช้งานโหมด Write-Ahead Logging (WAL) ซึ่งอนุญาตให้อ่านข้อมูลได้พร้อมกันหลายเธรด แต่การเขียนข้อมูลยังจำกัดให้มีผู้เขียนได้ทีละโพรเซส จึงต้องอาศัยกลไกการล็อกไฟล์ระดับเคอร์เนลเพื่อป้องกันปัญหาฐานข้อมูลถูกล็อกชั่วคราว
\end{enumerate}

% -----------------------------------------------------------
\vspace{-\baselineskip}
\section{ปัญหา อุปสรรค และแนวทางการแก้ไข}
% -----------------------------------------------------------

\begin{table}[H]
\centering
\footnotesize
\renewcommand{\arraystretch}{0.9}
\caption{ปัญหา อุปสรรค และแนวทางการแก้ไข}
\label{tab_problems}
\resizebox{\textwidth}{!}{%
\begin{tabular}{|c|L{6cm}|L{8cm}|}
\hline
\textbf{ลำดับ} & \textbf{ปัญหาที่พบ} & \textbf{แนวทางการแก้ไข} \\
\hline
1 & Cloudflare WAF บล็อกคำขอ HTTP ปกติ & ใช้ Playwright จำลองเบราว์เซอร์และเรนเดอร์ JavaScript \\ \hline
2 & บริการ LLM ภายนอกล่าช้าและมีโควตาจำกัด & ใช้ Deterministic Trending Engine พร้อมแคชและ Fallback \\ \hline
3 & ภาษาไทยไม่มีช่องว่างระหว่างคำและพาดหัวกำกวม & ใช้ PyThaiNLP ตัดคำ, thainer-v2 สกัดเอนทิตี และ WangchanBERTa สร้างเวกเตอร์ \\ \hline
4 & ข่าวเหตุการณ์เดียวกันแยกเป็นหลายกลุ่มข้ามวัน & ใช้ Persistent Story Identification, หน้าต่าง 72 ชั่วโมง, Two-Band และ Centroid EMA \\ \hline
5 & ข่าวประจำวันได้คะแนนนิยมสูงเกินจริง & ใช้ Routine News Dampening ลดคะแนนเหลือ 0.25 เท่า \\ \hline
6 & การเขียน JSON พร้อมกันทำให้ข้อมูลขัดแย้งและล็อกไฟล์ & ย้ายสู่ SQLite WAL และรองรับ MongoDB \\
\hline
\end{tabular}%
}
\end{table}

จากตารางที่~\ref{tab_problems} ปัญหาอุปสรรคทั้ง 6 ประการได้รับการแก้ไขจนสำเร็จลุล่วง ทำให้ระบบสามารถทำงานได้อย่างมีเสถียรภาพและรองรับการใช้งานต่อเนื่องได้อย่างมีประสิทธิภาพ

% -----------------------------------------------------------
\section{ข้อเสนอแนะสำหรับการพัฒนาในอนาคต}
% -----------------------------------------------------------

คณะผู้จัดทำมีข้อเสนอแนะในการพัฒนาโครงงานต่อยอดในอนาคต ดังนี้

\begin{enumerate}[noitemsep,topsep=2pt]
    \item \textbf{การพัฒนาระบบแจ้งเตือนแบบเรียลไทม์ (Real-Time Push Notification)}\\*
    พัฒนาระบบแจ้งเตือนแบบเรียลไทม์ผ่าน Web Push Notification หรือบริการส่งข้อความ เพื่อแจ้งเตือนผู้อ่านทันทีเมื่อมีพัฒนาการใหม่ในประเด็นข่าวที่ผู้อ่านได้บันทึกคั่นหน้าไว้
    \item \textbf{การวิจัยแบบจำลองตรวจจับอคติและจุดยืนของสื่อ (Media Bias and Stance Detection)}\\*
    วิจัยและพัฒนาโมเดลภาษาเพื่อจำแนกจุดยืนทางการเมืองและมุมมองการนำเสนอข่าวของแต่ละสำนักข่าว เพื่อช่วยให้ผู้อ่านได้รับทราบมุมมองที่หลากหลายและเป็นกลาง
\end{enumerate}
